\section{Conclusions}
\label{sec:conclusions}

Questo lavoro ha presentato un sistema di Federated Learning decentralizzato
peer-to-peer, in cui FedAvg è eseguito tramite aggregazione online e gossip
$k$-push senza alcun aggregatore centrale per l'intera durata
dell'addestramento, e in cui la centralizzazione è confinata alla sola
discovery dei client. Il sistema è stato valutato su FEMNIST (LEAF) con
partizionamento non-i.i.d.\ per scrittore, in 16 configurazioni tra esecuzione
locale e deployment reale su AWS multi-istanza, raggiungendo un'accuratezza
di validazione dell'87--91\% e mostrando che, nel regime testato, il costo di
comunicazione del protocollo gossip resta trascurabile rispetto al costo del
calcolo locale, anche su rete reale. L'analisi ha inoltre evidenziato, in
modo trasparente, i limiti della convergenza funzionale ottenuta e le
condizioni sperimentali (singola replica, un confondimento di configurazione)
sotto cui i risultati vanno interpretati.

Sviluppi futuri naturali includono: l'introduzione di una strategia di
aggregazione alternativa a FedAvg (ad es.\ FedProx) per mitigare
esplicitamente il client drift osservato; una politica di $k$ adattivo che
tenga conto del numero di peer ancora attivi; e la ripetizione delle
configurazioni con seed multipli per quantificare rigorosamente la varianza
run-to-run qui non misurata.

Due estensioni indirizzerebbero più a fondo la tolleranza ai guasti, oggi
solo parzialmente coperta (Sezione~\ref{sec:discussion}D). Per il rientro
dopo un crash, il checkpoint locale già salvato a ogni miglioramento della
validation loss potrebbe essere esteso con il numero di round: un worker
riavviato si ri-registrerebbe con lo stesso identificativo e riprenderebbe
da quel checkpoint invece che da zero, mentre il controllo di staleness già
presente assorbirebbe naturalmente il disallineamento di round rispetto agli
altri worker. Per l'ingresso live di un client nuovo, il discovery server
accetta già registrazioni in qualunque momento, ma oggi il nuovo worker
partirebbe da pesi casuali; aggiungendo al protocollo gRPC una RPC di
\emph{pull} accanto all'attuale \emph{push}, un worker entrante potrebbe
richiedere una sola volta i pesi correnti a un peer attivo come
inizializzazione, per poi partecipare ai round successivi come tutti gli
altri.
